\section{Discussion}
\label{sec:discussion}

\textbf{For agent designers.} The advice is conditional on the partner. Do not configure an agent to compete against an unknown partner: the rule pays only against a partner that cannot push back, pays by transfer, and is ruinous when met in kind. A \collaborating{} rule is the safe choice in the sense that it almost never lowers the joint payoff, but it pays its user only when the other side would not find the trades alone, that is with hidden interests, opaque options or a weaker default agent. Against a capable default agent no rule helps and the soft ones cost. When both sides can be configured, soft on both sides is the efficient cell once tokens cost anything.

\textbf{For the textbook.} The \tki{} prescriptions name real variables: every condition moved the map. But the direction is right in one case of six. What the model misses for LLM agents is that the strategies differ in what they need from the protocol (\collaborating{} needs moves and tokens) and that concessions are transfers unless both sides make them.

\textbf{For negotiation benchmarks.} What the names of the options reveal is a treatment. The same payoff structure produced a strategy effect when rendered abstractly and none when each side was told the other's needs, and the size of that difference depends on the model. A benchmark that reports a strategy or prompt effect on one rendering and one model has measured a point of the surface in \cref{fig:gap}.

\textbf{Limits.} Four models from two families and one cross-model pairing; newer models may close the remaining gaps. The strategies are rules we wrote, one operationalization among many, and \collaborating's trade clause is met in about half of its games. The payoff structure of the grounded instances (top, half, nothing) and the extractor's judgment are imposed. Games are one-shot, in English, with 20 to 40 instances per cell, so many entries of the condition map are directions; the price of tokens is swept rather than measured.

